\documentclass[11pt]{article}

% ===========================
% Style and packages
% ===========================

\usepackage{series}          % same house style as Papers I--III
\usepackage{amsmath,amssymb,amsthm}
\usepackage{geometry}
\usepackage{hyperref}
\usepackage{booktabs}

\geometry{margin=1in}

% ===========================
% Document
% ===========================

% --- Metadata ---
\title{Theta Closure in Modal Triplet Theory IV:
Conditional Gravity Scaling and Cosmological Cutoff Audit}
\author{Peter Nero}
\date{July 2026}

\begin{document}
\maketitle

\begin{abstract}
We re-evaluate the propagation of the gauge-profile geometry into gravity and
cosmology. Within the effective product ansatz of Papers~I--III, the updated
dimensionless internal-volume coefficient is
$\widehat V_{\mathrm{int}}=20.07064R_1^3$. Restoring the common internal length
$\ell_{\mathrm{int}}$ gives
$G_N^{-1}=20.07064\ell_{\mathrm{int}}^6R_1^3/G_{10}$; hence the gauge profile
does not determine Newton's constant without an absolute-scale and fundamental
gravity normalization. The former identification
$\Lambda_\Theta\sim5~\mathrm{TeV}$ is withdrawn with the obsolete gauge
crossing, so its numerical primordial-tensor bound is also withdrawn. We retain
the correctly normalized conditional relation
$r\leq 2\epsilon^2(\Lambda_\Theta/M_{\mathrm{Pl}})^2/(\pi^2A_s)$ when
$H\leq\epsilon\Lambda_\Theta$. The paper therefore supplies conditional
scaling laws and an assumption audit, not cross-sector numerical closure.
\end{abstract}


\section{Introduction}

Papers~I--III now establish a selected common-scheme gauge profile, a
calibrated leading geometric realization, and a conditional twistor
representation check. They do not select an absolute internal length, a
fundamental ten-dimensional gravitational coupling, or a cosmological
coherence cutoff. This distinction controls what can be transported into the
gravity and cosmology sectors.

This paper asks two scoped questions. First, what dimensionless internal-volume
coefficient follows if the effective round-$S^2$/nilmanifold product ansatz is
used? Second, what tensor inequality follows if a separately selected
coherence cutoff bounds the Hubble scale? The answers are conditional scaling
relations. They must not be promoted into predictions of $G_N$ or $r$ until
the missing absolute-scale and cutoff source theorems are supplied.


% ===========================
% ===========================
\section{Propagation of $\Theta$ to Newton's constant}
% ===========================

We compute the gravity scaling implied by the selected effective product
ansatz. The calculation introduces no retuning of the dimensionless profile,
but it retains the independent absolute length $\ell_{\mathrm{int}}$ and
fundamental coupling $G_{10}$.

\subsection{Gravitational coupling in Modal Triplet Theory}

In Modal Triplet Theory, the effective four--dimensional gravitational coupling
arises from dimensional reduction of the higher--dimensional coherent sector.
At leading order this takes the form
\begin{equation}
\frac{1}{G_N}
=
\frac{\mathrm{Vol}(X_{\mathrm{int}})}{G_{10}},
\label{eq:GN_reduction}
\end{equation}
where:
\begin{itemize}
\item $G_{10}$ is the fundamental gravitational coupling of the modal theory,
\item $X_{\mathrm{int}}$ is the internal coherent space.
\end{itemize}

Equation \eqref{eq:GN_reduction} is the standard dimensional-reduction
form assumed here. Its applicability to MTT requires the product metric,
Einstein-frame convention, and absence or control of warp/dilaton corrections.

\subsection{Internal volume in terms of $\Theta$}

In the $\Theta$--closure framework, the internal space is the product
\begin{equation}
X_{\mathrm{int}}
=
S^1_{\mathrm{cen}} \times \Sigma_2 \times \Sigma_3,
\end{equation}
where:
\begin{itemize}
\item $S^1_{\mathrm{cen}}$ is the central circle of radius $R_1$,
\item $\Sigma_2$ is the lens layer,
\item $\Sigma_3$ is the color layer $\Gamma\backslash\mathrm{Nil}_3$.
\end{itemize}

After factoring out a common physical length $\ell_{\mathrm{int}}$,
write the remaining radii and metric parameters as dimensionless quantities.
Under the unwarped product ansatz the physical internal volume factorizes as
\begin{equation}
\mathrm{Vol}(X_{\mathrm{int}})
=
\ell_{\mathrm{int}}^6(2\pi R_1)\;
\widehat{\mathrm{Area}}(\Sigma_2)\;
\widehat{\mathrm{Vol}}(\Sigma_3).
\label{eq:internal_volume_factorized}
\end{equation}

\subsection{Insertion of the selected gauge-profile geometry}

Paper~II's calibrated leading ansatz gives the dimensionless relations
\begin{align}
\widehat{\mathrm{Area}}(\Sigma_2)
&=4\pi(f_2R_{\mathrm{lens}})^2
=4\pi(0.2555137R_1),\\
\widehat{\mathrm{Vol}}(\Sigma_3)&=c=0.9948493R_1,
\qquad (a=b=1).
\end{align}
These are profile-realization data, not an independent gravity-sector
selection. Substitution gives
\begin{align}
\widehat V_{\mathrm{int}}
&=(2\pi R_1)\,[4\pi(0.2555137R_1)]\,(0.9948493R_1)\\
&=20.0706400R_1^3,
\end{align}
and therefore
\begin{equation}
\boxed{\mathrm{Vol}(X_{\mathrm{int}})
=20.0706400\,\ell_{\mathrm{int}}^6R_1^3.}
\label{eq:volume_result}
\end{equation}

\subsection{Conditional expression for $G_N$}

Under the dimensional-reduction assumptions,
\begin{equation}
\boxed{\frac{1}{G_N}
=\frac{20.0706400\,\ell_{\mathrm{int}}^6R_1^3}{G_{10}}.}
\label{eq:GN_theta}
\end{equation}
The gauge profile fixes only the displayed dimensionless shape coefficient
within the chosen ansatz. It does not fix $\ell_{\mathrm{int}}^6/G_{10}$, nor
does this paper derive the reduction formula from the selected MTT action.
Thus no numerical prediction of Newton's constant follows.


% ===========================
% ===========================
\section{Conditional cosmological cutoff relation}
% ===========================

The old gauge crossing did not select a physical coherence cutoff. In
particular, the former assignment $\Lambda_\Theta\sim5~\mathrm{TeV}$ and the
scan $[3,10]~\mathrm{TeV}$ are withdrawn. The current gauge matching point
$Q=M_t$ is a renormalization convention and must not be identified with
$\Lambda_\Theta$.

Let a future source theorem select a physical cutoff $\Lambda_\Theta$, and
suppose the relevant cosmological solution satisfies the quantitative
admissibility condition
\begin{equation}
H\leq\epsilon\Lambda_\Theta,
\qquad 0<\epsilon<1.
\label{eq:admissibility_H}
\end{equation}
For vacuum tensor fluctuations in standard slow-roll normalization,
\begin{equation}
P_t=\frac{2H^2}{\pi^2M_{\mathrm{Pl}}^2},
\qquad r=\frac{P_t}{A_s}.
\end{equation}
Consequently,
\begin{equation}
\boxed{r\leq
\frac{2\epsilon^2}{\pi^2A_s}
\left(\frac{\Lambda_\Theta}{M_{\mathrm{Pl}}}\right)^2.}
\label{eq:r_bound_general}
\end{equation}
This corrects the legacy formula, which omitted the factor
$2/(\pi^2A_s)$. Equation~\eqref{eq:r_bound_general} remains conditional on the
tensor-production model as well as on independently selected values of
$\Lambda_\Theta$ and $\epsilon$. It supplies no current numerical bound.


% ===========================
% ===========================
\section{Status and falsifiability}
% ===========================

The conditional relation~\eqref{eq:r_bound_general} can become falsifiable
only after MTT independently selects a cutoff, an admissibility margin, and a
cosmological production law. A measured tensor amplitude could then test that
joint hypothesis. Without those selections, neither detection nor
non-detection of primordial tensors presently falsifies the gauge-profile
closure program. The withdrawn $10^{-30}$--$10^{-29}$ interval must not be
quoted as an MTT prediction.


% ===========================
% ===========================
\section{Conclusion}
% ===========================

Within the selected effective product ansatz, the updated gauge profile fixes
the dimensionless internal-volume coefficient to
$20.0706400R_1^3$. Restoring dimensions shows explicitly why this is not yet a
prediction of Newton's constant: the independent combination
$\ell_{\mathrm{int}}^6/G_{10}$ remains.

The cosmological result is likewise a conditional scaling law. The obsolete
few-TeV gauge crossing cannot serve as a physical coherence cutoff, and the
legacy numerical tensor bound is withdrawn. A future closure theorem must
select $\Lambda_\Theta$, the quantitative margin $\epsilon$, and the applicable
cosmological state before Equation~\eqref{eq:r_bound_general} becomes a
numerical prediction. Paper~IV therefore documents cross-sector dependencies
and correct formulas; it does not establish gravity or cosmology closure.


% ===========================
% ===========================
\section*{References}
% ===========================

\begin{enumerate}

\item Particle Data Group (PDG), \emph{Review of Particle Physics},
for Standard Model input parameters and cosmological conventions.

\item M. E. Machacek and M. T. Vaughn,
``Two-loop renormalization group equations in a general quantum field theory,''
Nucl. Phys. B222 (1983) 83--103.

\item T. Kato, \emph{Perturbation Theory for Linear Operators},
Springer (1995).

\item P. Nero,
\emph{Modal Triplet Theory: Foundation},
MTT corpus.

\item P. Nero,
\emph{Modal Triplet Theory: Quantum Amplitudes from Modal Geometry},
MTT corpus.

\item P. Nero,
\emph{Direct Geometric Evaluation of Nonabelian Overlaps in Modal Triplet Theory},
Paper~II.

\item P. Nero,
\emph{Twistor--Action Matching of Gauge Overlaps in Modal Triplet Theory},
Paper~III.

\item R. Penrose and W. Rindler,
\emph{Spinors and Space-Time}, Vol. 2,
Cambridge University Press (for twistor geometry background).

\end{enumerate}
\begin{thebibliography}{99}

\bibitem{MTT-Admissibility}
P.~Nero,
\emph{Modal Triplet Theory: Admissibility, Encodings, and the Structure of Physical Description},
Zenodo preprint, January 2026.
\url{https://doi.org/10.5281/zenodo.18255621}

\bibitem{MTT-Foundation}
P.~Nero,
\emph{Modal Triplet Theory: Foundation},
Zenodo preprint, September 2025.
\url{https://doi.org/10.5281/zenodo.16949762}

\bibitem{MTT-FixedPoints}
P.~Nero,
\emph{Fixed Points I--VI: Complete Coherence Spine},
Zenodo preprints, August 2025.
\url{https://doi.org/10.5281/zenodo.16948748}

\bibitem{MTT-ProjectionAdmissibility}
P.~Nero,
\emph{The Projection--Admissibility Principle: Structural Constraints on Effective Physical Description},
Zenodo preprint, January 2026.
\url{https://doi.org/10.5281/zenodo.18255838}

\bibitem{MTT-Closure}
P.~Nero,
\emph{Closure and Inevitability in Modal Triplet Theory},
Zenodo preprint, January 2026.
\url{https://doi.org/10.5281/zenodo.18255510}

\bibitem{MTT-CoherenceCapacity}
P.~Nero,
\emph{Coherence Capacity as the Fundamental Resource of Effective Physics},
Zenodo preprint, January 2026.
\url{https://doi.org/10.5281/zenodo.18255905}

\bibitem{MTT-CoherenceDynamics}
P.~Nero,
\emph{Dynamics of Coherence Capacity: Transport, Concentration, and Exhaustion},
Zenodo preprint, January 2026.
\url{https://doi.org/10.5281/zenodo.18256048}

\bibitem{MTT-QM}
P.~Nero,
\emph{Modal Triplet Theory: From MTT to Quantum Mechanics},
Zenodo preprint, September 2025.
\url{https://doi.org/10.5281/zenodo.17074246}

\bibitem{MTT-QFT}
P.~Nero,
\emph{From MTT to Quantum Field Theory},
Zenodo preprint, 2025.
\url{https://doi.org/10.5281/zenodo.17068816}

\bibitem{MTT-GR}
P.~Nero,
\emph{Modal Triplet Theory: From MTT to General Relativity},
Zenodo preprint, October 2025.
\url{https://doi.org/10.5281/zenodo.16950597}

\bibitem{MTT-UVQG}
P.~Nero,
\emph{Modal Triplet Theory: From MTT to a UV-Finite, Unitary Quantum Gravity},
Zenodo preprint, 2025.
\url{https://doi.org/10.5281/zenodo.17077671}

\bibitem{MTT-Measurement}
P.~Nero,
\emph{Measurement as Disturbance and Stabilization in Modal Triplet Theory},
Zenodo preprint, 2025.
\url{https://doi.org/10.5281/zenodo.17177404}

\bibitem{MTT-Irreversibility}
P.~Nero,
\emph{Projection, Probability, and Irreversibility: Shadow Bridges Between Measurement, Black Holes, and Cosmology in Modal Triplet Theory},
Zenodo preprint, January 2026.
\url{https://doi.org/10.5281/zenodo.18256408}

\bibitem{MTT-Bell-Beables}
P.~Nero,
\emph{Modal Fixed Points, Bell’s Beables, and the Limits of Factorization},
Zenodo preprint, 2025.
\url{https://doi.org/10.5281/zenodo.17076300}

\bibitem{MTT-Bell-Temporal}
P.~Nero,
\emph{Temporal Bell Inequalities and Global Consistency in Modal Triplet Theory},
Zenodo preprint, August 2025.
\url{https://doi.org/10.5281/zenodo.18208884}

\bibitem{MTT-Stochastic}
P.~Nero,
\emph{From Modal Triplet Theory to Indivisible Stochastic Processes: A First-Principles, Fully Rigorous Derivation},
Zenodo preprint, January 2026.
\url{https://doi.org/10.5281/zenodo.18254862}

\end{thebibliography}

\end{document}
